\section{Virtualization and containerization}

\begin{frame}
	Virtualization and containerization.
\end{frame}

\begin{frame}
	\frametitle{Virtualization}
	\begin{columns}
		\begin{column}{0.5\textwidth}
			\begin{itemize}
				\item A host is a piece of hardware.
				\item A virtual machine takes pieces of host's CPU, memory, and
				      disc and simulates new hardware.
				\item A software called Hypervisor takes care of virtual
				      machines in a host.
				\item VirtualBox and VMware are popular virtualization
				      software.
			\end{itemize}
		\end{column}
		\begin{column}{0.5\textwidth}
			\begin{figure}
				\centering
				\includegraphics[scale=0.2]{logos/virtualbox.png}
				\includegraphics[scale=0.2]{logos/vmware.jpg}
			\end{figure}
		\end{column}
	\end{columns}
\end{frame}

\begin{frame}
	\frametitle{Containerization}
	\begin{columns}
		\begin{column}{0.5\textwidth}
			\begin{itemize}
				\item A host is a piece of hardware.
				\item A virtual machine takes pieces of host's CPU, memory, and
				      disc and simulates new hardware but it's able to share them
				      with other host's processes.
				\item Reproducible, lightweight environments for processes to
				      run.
				\item You're probably using them without knowing it.
				\item Docker and Podman are software for containerization.
			\end{itemize}
		\end{column}
		\begin{column}{0.5\textwidth}
			\begin{figure}
				\centering
				\includegraphics[scale=0.25]{logos/docker-logo.jpg}
				\includegraphics[scale=0.3]{logos/podman-logo.png}
			\end{figure}
		\end{column}
	\end{columns}
\end{frame}

\begin{frame}
	\frametitle{Docker}
	\begin{columns}
		\begin{column}{0.5\textwidth}
			\begin{itemize}
				\item An image is a template for creating containers.
				\item Images are immutable. Once created, they can't be
				      changed.
				\item A container is an instance of an image.
				\item Docker has a repository of images called
				      \href{https://hub.docker.com/}{dockerhub}.
			\end{itemize}
		\end{column}
		\begin{column}{0.5\textwidth}
			\begin{figure}
				\centering
				\includegraphics[scale=0.5]{logos/docker-logo.jpg}
			\end{figure}
		\end{column}
	\end{columns}
\end{frame}

\begin{frame}
	\frametitle{Rocker project}
	\begin{columns}
		\begin{column}{0.5\textwidth}
			\begin{itemize}
				\item Docker containers for the R
				      Environment.
				\item Docker images of r-base, rstudio, shiny, CUDA,
				      geospatial.
				\item Portable, sandboxed, transparent, community optimized,
				      versioned, extensible~\cite{boettiger2017}.
				\item Applications include
				      %development-debugging-testing,
				      data processing,
				      %deployment and continous delivery, create and
				      share computing environments (HPC), teaching, packaging
				      research reproducibility ~\cite{nust2020}.
				\item Rocker project web page \url{https://rocker-project.org}.
			\end{itemize}
		\end{column}
		\begin{column}{0.5\textwidth}
			\begin{figure}
				\centering
				\includegraphics[scale=0.25]{logos/rocker.png}
			\end{figure}
		\end{column}
	\end{columns}
\end{frame}

\begin{frame}
	\frametitle{Reproducible development environments}
	\begin{itemize}
		\item How do I get the same environment during developing and
		      processing?
	\end{itemize}
\end{frame}
